%!TEX program = lualatex
% SPDX-License-Identifier: CC-BY-SA-4.0
\documentclass[10pt]{article}

\usepackage{iftex}
\ifLuaTeX\else
  \PackageError{vintage-latex}{This example requires LuaLaTeX}
    {Run lualatex, not pdflatex or xelatex.}
\fi

\usepackage{fontspec}
\usepackage{xcolor}
\usepackage{geometry}
\usepackage{multicol}
\usepackage[tracking=true,protrusion=false,expansion=false]{microtype}

\defaultfontfeatures{Renderer=HarfBuzz}
\setmainfont{EB Garamond}[Numbers={OldStyle,Proportional}]
\definecolor{newsprint}{HTML}{F1E6CF}
\definecolor{printerink}{HTML}{211D18}
\pagecolor{newsprint}
\color{printerink}
\geometry{paperwidth=210mm,paperheight=210mm,
          top=17mm,bottom=17mm,left=17mm,right=17mm}
\pagestyle{empty}
\setlength{\columnsep}{6mm}
\setlength{\columnseprule}{.35pt}
\setlength{\parindent}{1em}
\setlength{\parskip}{0pt}

\newcommand{\articlehead}[1]{%
  \par\medskip
  \noindent\rule{\linewidth}{.55pt}\par
  \vspace{.25em}
  \begin{center}\large\textls[55]{\textsc{#1}}\end{center}
  \vspace{-.25em}
}

\begin{document}

\begin{center}
  {\fontsize{24.88}{26}\selectfont\textls[28]{THE SCIENTIFIC HERALD}}\par
  \vspace{.25em}
  \rule{\linewidth}{1.15pt}\par
  \vspace{.18em}
  \makebox[\linewidth]{%
    \footnotesize VOL. XII\hfill SATURDAY, 17 MAY 1879\hfill ONE PENNY}
  \vspace{.2em}
  \rule{\linewidth}{.35pt}
\end{center}

\begin{multicols}{2}

\articlehead{Electric Light in the Lecture Hall}

\noindent\textls[75]{\textsc{Last evening}} a compact arc lamp was exhibited
before the members of the Mechanics' Institute. Two carbon rods were advanced
by a clockwork movement, while the current was supplied by a small dynamo in
the adjoining room. The illumination was strong enough to project a spectrum
upon a screen at the far end of the hall.

The operator first showed the unequal consumption of the carbons, then reversed
the current so that the crater appeared upon the opposite rod. A white card,
held at several distances, made the rapid fall of brightness apparent without
requiring a photometer.

\begin{center}
  \vspace{.35em}
  \rule{.72\linewidth}{.35pt}\par
  \vspace{.45em}
  {\large\itshape ``A measurement is most useful when its conditions travel
  with it.''}\par
  \vspace{.45em}
  \rule{.72\linewidth}{.35pt}
  \vspace{.25em}
\end{center}

The demonstration was notable less for novelty than for order. Every change of
distance, polarity, and aperture was announced before the observation was
made. The record may therefore be repeated instead of merely admired.

\articlehead{Notes and Correspondence}

\noindent\textsc{A correspondent at Bath} reports that a magnetic needle,
suspended beneath a glass bell, retained its period after the air had been
partly withdrawn. Temperature was not recorded, and the observation should be
treated as provisional.

\smallskip
\noindent\textsc{Books received.}---A new table of logarithms; a practical
treatise on the construction of lenses; and the second part of a catalogue of
British mosses.

\columnbreak

\articlehead{The Transit Instrument}

\noindent\textls[75]{\textsc{The observatory}} has received a transit telescope
with a divided circle of twelve inches diameter. Its axis rests upon stone
piers separated from the surrounding floor, so that the motion of visitors is
not readily communicated to the instrument.

Before the first observations, the collimation, level, and azimuth errors were
determined separately. Repeating the same star upon both faces of the circle
then exposed an index error that a single reading would have concealed.

\articlehead{Meteorological Register}

\begin{center}
\small
\begin{tabular}{@{}lrrr@{}}
\textsc{Day} & \textsc{Bar.} & \textsc{Temp.} & \textsc{Rain}\\
\hline
Monday    & 29.84 & 54.2 & 0.00\\
Tuesday   & 29.71 & 51.8 & 0.13\\
Wednesday & 29.92 & 56.1 & 0.02\\
Thursday  & 30.06 & 58.7 & 0.00\\
Friday    & 30.11 & 60.3 & 0.00\\
\end{tabular}
\end{center}

\noindent\footnotesize Barometer reduced to 32 degrees Fahrenheit; rain in
inches. The ruled column divider, masthead, compact heads, and tabular register
are independent newspaper mechanisms rather than a single background texture.

\end{multicols}

\end{document}
